\documentclass[11pt]{article}
\usepackage{vinotheca}

\begin{document}

\vinothecaTitle{TasteRank Explorer}{Methods Primer}{A Guide to the Statistical and Network Methods}
\vinothecaAuthor{Jure Skarabot}{May 2026}
\vinothecaCompanions{Summary \textperiodcentered{} Technical Appendix \textperiodcentered{} Methods Primer \textperiodcentered{} Data Appendix \textperiodcentered{} Grape Reference}

\section*{Introduction}
The TasteRank Explorer uses a sequence of mathematical techniques from linear algebra, graph theory, and network science to convert wine tasting profiles into a similarity network, rank grape varieties by structural importance, and detect natural communities. This primer explains each method in the order it appears in the analytical pipeline: cosine similarity to measure resemblance between varieties, $k$-nearest-neighbour graph construction to build a network, eigenvector centrality to rank varieties within that network, PageRank as a robustness check, and modularity optimisation to detect communities. The treatment assumes comfort with basic linear algebra (matrix multiplication, eigenvalues) but not prior exposure to network science. Each section explains the intuition, states the formal definition, describes the algorithm used to compute it, and notes what the method reveals in the specific context of the TasteRank graph. For the full mathematical treatment with derivations, see the Technical Appendix.

\section{Cosine Similarity}

\textbf{What it does.} Cosine similarity measures how alike two grape varieties are based on the shape of their sensory profiles, independent of overall magnitude. Each variety is represented as a thirteen-dimensional vector (one dimension per sensory attribute: colour depth, aromatic intensity, floral character, and so on). Cosine similarity computes the cosine of the angle between two such vectors.

\textbf{Formal definition.} Given two profile vectors $\mathbf{p}$ and $\mathbf{q}$ in $\R^{13}$, the cosine similarity is the dot product of the vectors divided by the product of their magnitudes:
\[
\cos(\theta) = \frac{\mathbf{p} \cdot \mathbf{q}}{\norm{\mathbf{p}} \, \norm{\mathbf{q}}}.
\]
Because all sensory scores are non-negative integers (0--5), the result is always between 0 and 1. A score of 1 means the profiles are identical in shape (though they could differ in scale); 0 would mean the profiles share no common structure at all.

\textbf{Why cosine, not Euclidean distance?} Euclidean distance measures the absolute gap between vectors, so a variety with uniformly high scores would appear far from one with uniformly moderate scores even if their sensory shape is identical. Cosine similarity ignores this scale difference and focuses on the relative pattern of highs and lows across dimensions. This is more appropriate for tasting profiles, where the question is whether two varieties emphasise the same sensory attributes, not whether they score at the same absolute level.

\textbf{A practical subtlety.} The tannin dimension is structurally zero for most white grapes, creating a systematic difference between reds and whites. Cosine similarity handles this more gracefully than Euclidean distance: for whites, tannin contributes nothing to either the dot product or the norms, so it is effectively ignored rather than penalised. This is one reason the TasteRank network correctly identifies cross-boundary communities (C1: Light \& Aromatic) that include both reds and whites with similar non-tannin profiles.

\textbf{In the TasteRank graph.} The full $101 \times 101$ cosine similarity matrix has 5{,}050 unique pairwise values. The mean similarity across all connected pairs is 0.983, with a narrow range of $[0.909, 1.000]$. This tight distribution reflects the fact that wine grape profiles share a common baseline shape --- the differentiation between varieties lies in subtle but structurally meaningful differences across specific dimensions.

\section{$k$-Nearest-Neighbour Graph Construction}

\textbf{What it does.} The similarity matrix is dense --- every variety has a non-zero similarity to every other. A $k$-nearest-neighbour ($k$NN) graph sparsifies this by keeping only the strongest relationships: each variety is connected by an edge to its $k$ most similar neighbours. The result is a graph where nodes are grape varieties and weighted edges represent strong sensory resemblance.

\textbf{How it works.} For each variety, the algorithm sorts all other varieties by descending cosine similarity and retains the top $k = 5$. This produces a directed graph (variety A may include B in its top 5, but B may not include A). The graph is then symmetrised: an undirected edge is placed between A and B if either A lists B or B lists A among their five nearest neighbours. The edge weight is the cosine similarity between the pair.

\textbf{Why $k = 5$?} The choice of $k$ balances two competing goals. A small $k$ produces a sparse graph that highlights only the very strongest similarities, at the risk of fragmenting the network into disconnected components. A large $k$ produces a dense graph where every variety is connected to many others, drowning out the most meaningful relationships. At $k = 5$, the TasteRank graph is connected (no isolated nodes), sparse enough to reveal meaningful structure, and produces a clear community partition.

\textbf{Symmetrisation and degree inflation.} Because the graph is symmetrised by union (not intersection), the effective degree of each node is typically higher than $k$. The average degree in the TasteRank graph is 6.75, reflecting this inflation above the nominal $k = 5$.

\textbf{In the TasteRank graph.} The $k$NN construction produces 341 edges among 101 nodes, with a graph density of 0.068 (6.8\% of all possible edges). This sparse, weighted graph is the input to all subsequent analyses --- centrality, PageRank, and community detection all operate on this structure.

\section{Eigenvector Centrality (TasteRank)}

\textbf{What it does.} Eigenvector centrality assigns a score to each node in a network that reflects not just how many connections it has, but how important its connections are. A variety earns a high TasteRank score by being similar to other varieties that are themselves highly central --- a recursive definition that rewards membership in densely connected, mutually reinforcing clusters.

\subsection{The eigenvalue problem}
The weighted adjacency matrix $W$ is a $101 \times 101$ symmetric matrix where entry $W_{ij}$ is the cosine similarity between varieties $i$ and $j$ if they share an edge, and zero otherwise. Being real and symmetric, $W$ has 101 real eigenvalues $\lambda_1 \geq \lambda_2 \geq \cdots \geq \lambda_{101}$ and corresponding orthogonal eigenvectors $\mathbf{x}_1, \mathbf{x}_2, \ldots, \mathbf{x}_{101}$, satisfying
\[
W \mathbf{x}_k = \lambda_k \mathbf{x}_k \qquad \text{for each } k.
\]
Eigenvector centrality defines the importance of each variety as its component in the leading eigenvector $\mathbf{x}_1$ (the eigenvector associated with the largest eigenvalue $\lambda_1$). The TasteRank score of variety $i$ is simply $x_{1,i}$, normalised so the eigenvector has unit length.

\subsection{The recursive interpretation}
Written component-wise, the eigenvector equation says that
\[
x_i = \frac{1}{\lambda_1} \sum_j W_{ij}\, x_j.
\]
In words: each variety's centrality is proportional to the weighted sum of its neighbours' centralities, which are defined by their neighbours' centralities, and so on to infinite depth. This circularity is not a defect --- it is the defining feature. The eigenvector is the unique self-consistent solution to this infinite chain of mutual dependencies. This is directly analogous to Google's PageRank: there is no closed-form formula for any individual page's rank.

\subsection{Why the largest eigenvalue dominates}
Consider starting with any initial vector of scores $\mathbf{z}_0$ and repeatedly multiplying by $W$. After $t$ iterations, $\mathbf{z}_t = W^t \mathbf{z}_0$. Decomposing $\mathbf{z}_0$ in the eigenbasis, the contribution of each eigenvector is scaled by $\lambda_k^t$. Since $\lambda_1$ is strictly the largest eigenvalue (guaranteed by the Perron--Frobenius theorem for connected, non-negative matrices), the ratio $|\lambda_2 / \lambda_1|^t$ converges to zero exponentially. After enough iterations, only the leading eigenvector survives. This is why $\lambda_1$ alone determines the centrality ranking.

\subsection{The Perron--Frobenius theorem}
The theoretical guarantee behind eigenvector centrality is the Perron--Frobenius theorem, which states that for a connected graph with non-negative edge weights: (1) the largest eigenvalue $\lambda_1$ is real, positive, and strictly greater than the magnitude of all other eigenvalues; (2) the corresponding eigenvector $\mathbf{x}_1$ has all positive components; (3) $\mathbf{x}_1$ is unique up to normalisation. Condition (2) is essential --- it ensures that centrality scores are all positive and thus interpretable as a ranking. Condition (3) ensures that the ranking is unique and does not depend on the choice of algorithm or starting point.

\subsection{The spectral gap}
The spectral gap $|\lambda_2 / \lambda_1|$ controls both the convergence speed of the computation and the structural interpretation. A small ratio means fast convergence and a clear, dominant centrality structure. A ratio near 1 would indicate competing centres of influence with no clear hierarchy. For the TasteRank graph, the spectral gap is approximately 0.72, meaning each iteration of power iteration reduces the second-order contamination by approximately 28\%. After 25 iterations, the residual error is below $10^{-3}$; after 50, below $10^{-7}$.

\subsection{Power iteration}
The standard algorithm for computing the leading eigenvector is power iteration. Starting from a uniform vector, the algorithm repeatedly multiplies by $W$ and normalises:
\[
\mathbf{x}_{t+1} = \frac{W \mathbf{x}_t}{\norm{W \mathbf{x}_t}}.
\]
This converges to $\mathbf{x}_1$ because the repeated multiplication amplifies the leading eigenvector component and suppresses all others. The eigenvalue is recovered as the Rayleigh quotient $\lambda_1 = \mathbf{x}^\top W \mathbf{x}$. In the TasteRank graph, convergence to tolerance $10^{-6}$ is achieved in 35--45 iterations.

\textbf{In the TasteRank graph.} Sagrantino ranks first (TasteRank $= 0.3020$) because it is connected to Nero d'Avola, Lagrein, Montepulciano, Petite Sirah, and Petit Verdot --- all top-10 varieties themselves embedded in the dense Mediterranean red cluster (C0). The recursive amplification through this mutually reinforcing core pushes it to the top. Conversely, Riesling ranks low because its neighbours (Albari\~no, Furmint, Sauvignon Blanc) are themselves peripheral, and the recursive logic works in reverse.

\section{PageRank}

\textbf{What it does.} PageRank, developed by Brin and Page (1998) for ranking web pages, is a modified form of eigenvector centrality that includes a teleportation mechanism to prevent excessive concentration of importance within a single dense cluster. In TasteRank, it serves as a robustness check on the eigenvector centrality ranking.

\subsection{The random surfer model}
A random taster navigates the grape similarity network. At each step, the taster either (a) with probability $\alpha$ follows an edge to a similar variety, choosing among neighbours proportionally to cosine similarity weight, or (b) with probability $1 - \alpha$ teleports to a completely random variety. The PageRank of each variety is the long-run fraction of time the random taster spends there. Formally, the transition is governed by the matrix
\[
G = \alpha M + (1 - \alpha) \frac{1}{n} J,
\]
where $M$ is the row-normalised adjacency matrix and $J$ is the all-ones matrix. The PageRank vector $\boldsymbol\pi$ is the stationary distribution: $G^\top \boldsymbol\pi = \boldsymbol\pi$. With the standard damping factor $\alpha = 0.85$, the random taster follows edges 85\% of the time and teleports 15\% of the time.

\subsection{How it differs from eigenvector centrality}
The teleportation mechanism ensures that every variety receives a baseline importance floor of $(1 - \alpha)/n$, regardless of its network position. This prevents the extreme concentration that can occur with pure eigenvector centrality in networks with dense subgraphs. PageRank also rewards bridge varieties --- nodes connecting multiple communities --- because the random walk must pass through them when transitioning between clusters. The two measures agree strongly overall (Spearman rank correlation $\rho \approx 0.92$) but diverge in informative ways. Varieties deep inside a dense cluster have higher eigenvector centrality but slightly lower PageRank. Bridge varieties show the opposite pattern. The most striking example is Sangiovese: TasteRank rank 40 vs.\ PageRank rank $\sim 8$.

\textbf{Interpreting divergences.} When PageRank $\gg$ TasteRank for a variety, it is acting as a structural bridge between communities. When TasteRank $\gg$ PageRank, the variety is embedded deep in a dense cluster core. The gap between the two measures is itself a diagnostic for network position.

\section{Modularity and Community Detection}

\textbf{What it does.} Community detection partitions the graph into groups where within-group connectivity is stronger than would be expected by chance. The result is a set of clusters that reflect natural groupings in the data --- in this case, families of grape varieties with mutually similar sensory profiles.

\subsection{The modularity function}
The modularity function $Q$ measures the quality of a partition. For each pair of nodes in the same community, it compares the actual edge weight to the expected weight under a null model (the configuration model, which randomly rewires edges while preserving each node's total connectivity). Formally:
\[
Q = \frac{1}{2m} \sum_{i,j} \left[ W_{ij} - \frac{s_i s_j}{2m} \right] \delta(c_i, c_j),
\]
where $m$ is total edge weight, $s_i$ is the weighted degree of node $i$, $c_i$ is its community assignment, and $\delta$ is the Kronecker delta. A partition with $Q > 0.3$ is generally considered to have significant community structure.

\subsection{The null model}
The null model deserves emphasis because it defines what ``expected by chance'' means. The configuration model preserves the degree sequence --- each node retains its total edge weight --- but randomises which specific nodes are connected. Under this model, the expected edge weight between nodes $i$ and $j$ is $s_i s_j / 2m$. Nodes with high total connectivity are expected to share more edge weight simply because they are well-connected, not because they are in the same community. Modularity measures the excess above this baseline, isolating signal from noise.

\subsection{The greedy algorithm}
The Clauset--Newman--Moore algorithm is an agglomerative method that starts with every node in its own singleton community and iteratively merges the pair of communities whose merger produces the largest increase in $Q$. The sequence of merges defines a dendrogram, and the partition with the highest $Q$ across the entire merge history is selected as the final community structure. This greedy approach does not guarantee the globally optimal partition (modularity optimisation is NP-hard), but it is fast and produces results competitive with more expensive methods. For the TasteRank graph, it was validated against the Louvain method and spectral bisection, all three producing partitions with similar membership and comparable modularity.

\subsection{Connection to the eigenvalue spectrum}
There is a deep connection between community structure and the eigenvalues of the modularity matrix $B$ (where $B_{ij} = W_{ij} - s_i s_j / 2m$). The number of positive eigenvalues of $B$ provides an upper bound on the number of meaningful communities. For the TasteRank graph, the first six eigenvalues of $B$ are positive and clearly separated from the bulk spectrum, consistent with the six-community partition.

\textbf{In the TasteRank graph.} Six communities are detected with modularity $Q \approx 0.41$, indicating strong structure. The partition tracks sensory logic rather than geography: Nerello Mascalese clusters with Pinot Noir in C2 (mid-weight structured reds), not with its Sicilian neighbour Nero d'Avola in C0 (Mediterranean reds). Riesling and Assyrtiko join the cross-boundary C1 cluster rather than the mineral-white C4, because their extreme acidity and complexity scores align them with light aromatic varieties.

\section{The Pipeline in Summary}

The five methods form a coherent analytical pipeline, each building on the output of the previous step. (1) \emph{Cosine similarity} converts 101 sensory profile vectors into a $101 \times 101$ similarity matrix, measuring shape-based resemblance between every pair. (2) \emph{$k$NN graph construction} sparsifies the dense similarity matrix into a weighted network of 341 edges, retaining only the strongest relationships ($k = 5$ per variety, symmetrised). (3) \emph{Eigenvector centrality} ranks varieties by their recursive structural importance in this network --- this is the TasteRank score. (4) \emph{PageRank} provides a robustness check with a damping mechanism that prevents excessive concentration; divergences identify bridge varieties and cluster cores. (5) \emph{Modularity optimisation} partitions the network into six communities of varieties with mutually similar profiles.

The mathematical thread connecting these steps is the weighted adjacency matrix $W$, which encodes the network's structure. Eigenvector centrality uses the leading eigenvector of $W$; PageRank uses a damped, row-normalised version of $W$; and modularity detection uses a null-model corrected version of $W$. The same matrix, analysed from three different perspectives, yields the TasteRank ranking, the robustness diagnostic, and the community partition.

\section*{References}

\begin{itemize}[leftmargin=2em,itemsep=0.2em]
\item Bonacich, P. (1987). Power and centrality: A family of measures. \emph{American Journal of Sociology}, 92(5), 1170--1182.
\item Brin, S. \& Page, L. (1998). The anatomy of a large-scale hypertextual web search engine. \emph{Computer Networks and ISDN Systems}, 30(1--7), 107--117.
\item Clauset, A., Newman, M. E. J. \& Moore, C. (2004). Finding community structure in very large networks. \emph{Physical Review E}, 70(6), 066111.
\item Fruchterman, T. M. J. \& Reingold, E. M. (1991). Graph drawing by force-directed placement. \emph{Software: Practice and Experience}, 21(11), 1129--1164.
\item Meyer, C. D. (2000). \emph{Matrix Analysis and Applied Linear Algebra}. SIAM.
\item Newman, M. E. J. (2006). Modularity and community structure in networks. \emph{Proceedings of the National Academy of Sciences}, 103(23), 8577--8582.
\item Newman, M. E. J. (2010). \emph{Networks: An Introduction}. Oxford University Press.
\end{itemize}

\end{document}
